% Αρχιτεκτονική του συστήματος σε επίπεδα (περιεχόμενο σχήματος· το περιβάλλον
% figure και η λεζάντα ορίζονται στο Κεφάλαιο 5). Τα ονόματα αρθρωμάτων
% αντιστοιχούν στο πραγματικό δέντρο πηγαίου κώδικα: src/rubik_optimal/
% (cli.py, oracle.py, verify.py, solvers/, tables/), native/ και results/.
\begin{tikzpicture}[
  font=\footnotesize,
  band/.style={draw=black!60, rounded corners=3pt, fill=black!4},
  bandtitle/.style={anchor=north west, font=\scriptsize\bfseries, text=black!70, inner sep=2pt},
  comp/.style={draw=black!75, rounded corners=1.5pt, fill=white, align=center,
               font=\scriptsize, inner sep=3pt, minimum height=1.05cm},
  flow/.style={-{Stealth[length=2.4mm]}, semithick},
  resflow/.style={-{Stealth[length=2.4mm]}, semithick, black!55},
  lbl/.style={font=\tiny, align=center}
]

% Επίπεδο 1: Διεπαφή γραμμής εντολών
\draw[band] (0,14.3) rectangle (12.4,15.9);
\node[bandtitle] at (0.12,15.82) {Διεπαφή γραμμής εντολών (CLI)};
\node[comp, minimum width=10.8cm] at (6.2,14.85)
  {{\ttfamily rubik-optimal} ({\ttfamily src/rubik\_optimal/cli.py}): οκτώ υποεντολές\\
   {\tiny\ttfamily solve · distance · oracle · benchmark · verify · scramble · facelets · tables}};

% Επίπεδο 2: Επιλυτές Python
\draw[band] (0,11.3) rectangle (12.4,13.5);
\node[bandtitle] at (0.12,13.42) {Επιλυτές Python ({\ttfamily rubik\_optimal.solvers})};
\node[comp, minimum width=2.85cm] at (1.65,12.2)
  {Korf / IDA*\\ {\tiny\ttfamily korf.py}};
\node[comp, minimum width=2.85cm] at (4.68,12.2)
  {Kociemba\\ {\tiny\ttfamily kociemba.py}\\ {\tiny\ttfamily kociemba\_optimal.py}\\
   {\tiny προσαρμογέας {\ttfamily kociemba} (PyPI)}};
\node[comp, minimum width=2.85cm] at (7.71,12.2)
  {Thistlethwaite\\ {\tiny\ttfamily thistlethwaite.py}};
\node[comp, minimum width=2.85cm] at (10.74,12.2)
  {Κύβος 2×2×2 (\eng{Pocket})\\ {\tiny\ttfamily pocket\_cube.py · pocket/}};

% Στρώμα μαντείου (γέφυρα επιπέδων 2-3)
\node[comp, fill=black!8, minimum width=7.4cm] at (8.4,10.05)
  {Στρώμα μαντείου: {\ttfamily oracle.py}: {\ttfamily FastOptimalOracle}\\
   {\tiny παραμένουσα διεργασία (\eng{resident}) και λειτουργία δέσμης (\eng{batch})}};

% Επίπεδο 3: Εγγενείς μηχανές C/C++
\draw[band] (1.6,6.85) rectangle (12.4,8.85);
\node[bandtitle] at (1.72,8.77) {Εγγενείς μηχανές C/C++ ({\ttfamily native/})};
\node[comp, minimum width=2.5cm] at (3.01,7.6)
  {Εγγενής IDA*\\ {\tiny\ttfamily optimal\_solver}};
\node[comp, minimum width=2.5cm] at (5.67,7.6)
  {Παραγωγή PDB\\ {\tiny\ttfamily corner\_pdb}\\ {\tiny\ttfamily edge\_pdb}};
\node[comp, minimum width=2.5cm] at (8.33,7.6)
  {Απόδειξη φάσης 2\\ {\tiny\ttfamily kociemba\_phase2\_probe}};
\node[comp, minimum width=2.5cm] at (10.99,7.6)
  {Μηχανή H48\\ {\tiny\ttfamily h48\_backend}\\ {\tiny ενσωματωμένο {\ttfamily nissy-core}}};

% Επίπεδο 4: Πίνακες
\draw[band] (0,4.05) rectangle (12.4,6.05);
\node[bandtitle] at (0.12,5.97) {Πίνακες ({\ttfamily data/generated/})};
\node[comp, minimum width=3.7cm] at (2.17,4.8)
  {PDB γωνιών + 8 PDB έξι ακμών\\
   {\tiny\ttfamily *\_corner\_state\_pdb.bin}\\ {\tiny\ttfamily *\_edge\_subset\_*\_pdb.bin}};
\node[comp, minimum width=3.7cm] at (6.19,4.8)
  {Πίνακες αποστάσεων\\ συμπλόκων\\
   {\tiny Thistlethwaite · φάση 1 Kociemba}};
\node[comp, minimum width=3.7cm] at (10.21,4.8)
  {Πίνακας αποκοπής H48\\ {\tiny\ttfamily h48h7.bin}};

% Επίπεδο 5: Επαλήθευση και αποτελέσματα
\draw[band] (0,1.05) rectangle (12.4,3.25);
\node[bandtitle] at (0.12,3.17) {Επαλήθευση και αποτελέσματα};
\node[comp, minimum width=3.7cm] at (2.17,2.0)
  {Ελεγκτής λύσεων\\ {\tiny\ttfamily verify.py}\\ {\tiny\ttfamily scripts/verify\_results.py}};
\node[comp, minimum width=3.7cm] at (6.19,2.0)
  {Αρχεία αποτελεσμάτων\\ {\tiny\ttfamily results/raw $\to$ results/processed}};
\node[comp, minimum width=3.7cm] at (10.21,2.0)
  {Εξωτερικός διασταυρωτικός\\ έλεγχος: Nissy 2.0.8\\ {\tiny\ttfamily solvers/nissy\_external.py}};

% Ροές κλήσεων (προς τα κάτω)
\draw[flow] (6.2,14.3) -- (6.2,13.5)
  node[lbl, right=2pt, midway] {κλήσεις επιλυτών και μαντείου};
\draw[flow] (8.4,11.3) -- (8.4,10.5)
  node[lbl, right=2pt, midway] {ερωτήματα ακριβούς απόστασης /\\ βέλτιστης λύσης};
\draw[flow] (8.4,9.6) -- (8.4,8.85)
  node[lbl, right=2pt, midway] {παραμένουσα εγγενής διεργασία};
\draw[flow] (3.0,11.3) -- (3.0,8.85)
  node[lbl, left=2pt, midway, align=right]
  {περιτυλίγματα\\ {\ttfamily optimal\_native.py}\\ {\ttfamily h48\_native.py}};
\draw[flow] (0.7,11.3) -- (0.7,6.05);
\node[lbl, rotate=90] at (0.42,8.68) {απευθείας φόρτωση πινάκων};
\draw[flow] (6.19,6.85) -- (6.19,6.05)
  node[lbl, right=2pt, midway] {παραγωγή και φόρτωση πινάκων};

% Ροή αποτελεσμάτων (δεξιός διάδρομος)
\draw[resflow] (12.4,12.2) -- (13.25,12.2) -- (13.25,2.0) -- (12.4,2.0);
\draw[semithick, black!55] (12.4,7.6) -- (13.25,7.6);
\node[lbl, rotate=90, text=black!60] at (13.6,7.1)
  {ροή εγγραφών αποτελεσμάτων (JSONL)};

\end{tikzpicture}
